% ================================================
% 文件: 02_有线广播.tex
% 章节: 第15章 有线广播
% 所属部分: 系统与设备
% 版本: v1.0
% ================================================
% 本章目录结构（树状）:
% \chapter{有线广播}
% ├── \section{有线广播概述}
% ├── \section{有线广播系统组成}
% ├── \section{有线广播传输技术}
% ├── \section{有线广播设备}
% └── \section{有线广播设计与安装}
% ================================================

\chapter{有线广播}
\label{chap:wired_broadcast}

\section{有线广播概述}
\label{sec:wired_broadcast_intro}

有线广播是通过有线传输介质（如电缆、双绞线、光纤）传输音频信号的
广播系统。与无线广播相比，有线广播具有信号稳定、抗干扰能力强、音
质好、不易被截获等优点，广泛用于需要可控覆盖与高可靠播发的场合。
它不占用无线频谱资源，也不受电离层与天气影响，因而在一些关键场景
中具有不可替代的地位。

有线广播的应用场景十分广泛：

\begin{itemize}
    \item \textbf{公共广播}：学校、商场、车站
、机场等公共场所的背景音乐与业务广播。
    \item \textbf{应急广播}：火灾报警、紧急疏
散通知、自然灾害预警，要求高可靠与强制播发。
    \item \textbf{校园广播}：教学广播、作息铃
声、考试指令，强调分区与定时。
    \item \textbf{楼宇广播}：背景音乐、寻呼通
知，常与消防系统联动。
\end{itemize}

从技术演进看，有线广播正由传统的模拟定压/定阻系统，逐步向"IP
 化、网络化、智能化"发展：以局域网或专用传输网承载数字音频，借
助软件实现灵活分区、远程播控与多级联动。但在农村、厂矿与既有建筑
中，模拟定压系统因其简单、廉价、抗网络故障，仍大量存在并构成应急
广播的物理基础。无论模拟还是数字，衡量广播质量的核心仍是频率响应
、失真度与信噪比三大电声指标。

有线广播的设计目标是在覆盖区内提供清晰、均匀且足够响度的声音，同
时做到分区可控、应急优先与运行可靠。为此，系统须兼顾声学（扬声器
布点、声压级）与电学（功率、阻抗、压降）两方面：声学决定"听得清
、听得匀"，电学决定"供得上、不烧机"。在公共与应急领域，相关标
准对声压级、语言传输指数（STI）与备用电源续航提出了明确要求，
设计时应作为底线而非上限来对待。

在防灾减灾体系中，农村应急广播承担着"最后一公里"的信息触达使命
：县级播控平台统一制作与下发预警，经有线（或结合无线/卫星）链路
送达乡镇与行政村的大喇叭、音柱，并具备断网可本地播发的兜底能力。
这类系统强调无人值守、远程巡检与多通道备份，常与气象、地震、防汛
等部门联动，在灾害发生时第一时间发布避险指令。其可靠性不取决于单
一技术先进与否，而取决于"传输冗余 + 终端可靠 + 供电备份"
的整体韧性。

\section{有线广播系统组成}
\label{sec:wired_broadcast_components}

有线广播系统通常由音源、处理、放大、传输与终端五部分构成：

\begin{itemize}
    \item \textbf{音源设备}：麦克风、CD 播
放器、MP3/数字播放器、调音台、广播控制台，提供话音与节目信号
。
    \item \textbf{功率放大器}：将音频信号放大
到足以驱动远距离扬声器的功率电平，是有线广播的功率核心。
    \item \textbf{传输线路}：音频电缆、双绞线
、同轴缆或光纤，负责将功放输出送达各分区。
    \item \textbf{终端设备}：吸顶/壁挂扬声器
、音柱、号角喇叭，将电信号还原为声能。
    \item \textbf{控制设备}：分区控制器、定时
器、矩阵切换器、紧急呼叫器，实现分区寻呼与优先级播发。
\end{itemize}

系统工作时，音源信号经调音/前置放大与必要的均衡处理后送入功率放
大器，再沿传输线路分配到各分区扬声器。为保证重要信息优先播发，控
制设备通常设置优先级：消防与应急信号可强切背景音乐并解锁指定分区
。对于大型系统，常采用"总控机房—分控/功放站—末端扬声器"三级
架构，既便于集中管理，也降低单点故障影响范围。传输线路的阻抗与压
降设计是系统成败的关键，需在下一节结合定压/定阻方式详细讨论。

在集中式系统中，常设前置放大器（话放/线路放大）对音源进行电平调
整、均衡与混音，再送功率放大器；而"线路放大器"用于远距离信号传
输前提升电平、降低线路噪声比，与末端的定压功放分工明确。分区控制
则通过分区器或矩阵将同一音源或不同音源送达指定区域：平时分区播放
背景音乐，紧急时由优先级最高的消防/应急信号强切并解锁相关分区，
实现"局部寻呼、全局强播"的灵活策略。分区划分应与建筑防火分区、
功能用房一致，既便于管理也满足疏散引导的需要。

\section{有线广播传输技术}
\label{sec:wired_broadcast_transmission}

有线广播主要有定压、定阻与数字（IP/光纤）三类传输方式，其适用
场合差异明显。

\begin{itemize}
    \item \textbf{定压传输}：功放输出标称 $
70\ \mathrm{V}$ 或 $100\ \mathrm{
V}$ 高压小电流信号，适合远距离、多扬声器并联的集中系统。
    \item \textbf{定阻传输}：功放输出低电压、
低内阻，直接驱动定阻扬声器，适合短距离、高保真（如厅堂扩声）。
    \item \textbf{数字传输}：采用 IP 网络
或专用数字总线传输，支持多通道、远程控制与双向状态回传。
    \item \textbf{光纤传输}：以光信号承载音频
，适合长距离、强干扰或需电气隔离的场合。
\end{itemize}

定压传输是有线广播的主流方案，其优点在于线路损耗小、便于多个扬声
器并联使用。在定压系统中，每个扬声器经"线间变压器"接入高压母线
，变压器既实现高压/低压变换，又完成阻抗匹配。设功放标称输出电压
为 $U_0$，扬声器额定功率为 $P_s$、标称阻抗为 $R_
s$，则变压器电压变比 $n$ 与阻抗变换满足
\begin{equation}
n = \frac{U_0}{U_s}, \qquad Z_{\mathrm{in}} = n^2 R_s = \left(\frac{U_0}{U_s}\right)^2 R_s,
\end{equation}
其中 $U_s$ 为扬声器端电压，$Z_{\mathrm{in}
}$ 为折算到高压侧的等效阻抗。多只扬声器并联时，系统总功率近似
为各只功率之和，母线电流随之增加。线路压降 $\Delta U$
 与导线电阻 $R_L$、母线电流 $I$ 有关：
\begin{equation}
\Delta U = I R_L, \qquad P_{\mathrm{loss}} = I^2 R_L,
\end{equation}
故远距离传输应加大线径以降低 $R_L$，把线路损耗控制在总功率
的一成以内为典型要求。

\begin{table}[htbp]
\centering
\caption{定压与定阻传输方式对照}
\label{tab:wired_mode}
\begin{tabular}{|p{2.6cm}|p{3.6cm}|p{4.0cm}|p{3.0cm}|}
\hline
\textbf{方式} & \textbf{标称电压} & \textbf{适用距离} & \textbf{主要优点} \\
\hline
定压 $70\ \mathrm{V}$ & $70\ \mathrm{V}$ & 中远距离 & 并联方便、压降小 \\
\hline
定压 $100\ \mathrm{V}$ & $100\ \mathrm{V}$ & 远距离 & 损耗更低、更常用 \\
\hline
定阻 & 低电压 & 近距离 & 高保真、音质好 \\
\hline
数字/IP & 网络电平 & 任意（网内） & 灵活分区、可网管 \\
\hline
\end{tabular}
\end{table}

定压系统的线路压降可用前述 $\Delta U = I R_L$
 估算。例如母线电压 $100\ \mathrm{V}$、总功率
 $500\ \mathrm{W}$ 时母线电流 $I=5\ \
mathrm{A}$；若馈线往返总阻 $R_L=0.4\ \Om
ega$，则压降 $\Delta U=2\ \mathrm{V}
$（约 $2\%$），线路损耗仅 $10\ \mathrm{W}
$，处于可接受范围。若距离翻倍使 $R_L$ 增至 $0.8\ 
\Omega$，压降升至 $4\ \mathrm{V}$、损耗达
 $20\ \mathrm{W}$，此时应加粗线径或提高标称电压
（如改 $70\ \mathrm{V}$ 为 $100\ \ma
thrm{V}$）以降低电流。定阻系统则须保证功放标称阻抗与扬声
器总阻抗匹配，避免欠载导致功放过压、或过载导致功放过热。

数字（IP）广播以局域网或专用传输网承载编码音频流，每个终端为带
解码与功放的网络音箱，优势是可灵活分区、集中网管与双向状态回传，
且易于与消防、安防平台集成；其代价是对网络带宽、时钟同步与供电（
常需 PoE）有要求。光纤传输则用于强电磁干扰、长距离或需电气隔
离的骨干链路，将音频调制为光信号后几乎不受地环路与雷击影响。

\section{有线广播设备}
\label{sec:wired_broadcast_equipment}

主要设备及其选型要点如下：

\begin{itemize}
    \item \textbf{功率放大器}：定压功放（输出
带线间变压器抽头）、定阻功放（用于高保真扩声）、数字功放（效率高
、体积小）。容量应按扬声器总功率并预留 $20\%\sim 30
\%$ 余量。
    \item \textbf{扬声器}：吸顶扬声器、壁挂扬
声器、音柱、号角喇叭，按声压级、指向性与环境防护等级选择。
    \item \textbf{音源设备}：专业麦克风、广播
控制台、数字音频播放器、调谐器，要求低噪声与稳定输出。
    \item \textbf{控制设备}：分区器、矩阵切换
器、紧急呼叫器、定时器，负责寻呼、强切与优先级管理。
\end{itemize}

定压功放常在输出变压器上设置多个功率抽头（如 $1/4$、$1/
2$、全功率），以适应不同分区负荷。扬声器侧线间变压器亦设有功率
档（如 $3\ \mathrm{W}$、$6\ \mathrm{
W}$、$10\ \mathrm{W}$、$15\ \mathr
m{W}$），安装时按设计功率选择抽头，确保各扬声器端电压一致、
声压均匀。对于号角等高灵敏度扬声器，小功率即可覆盖大区域，适合车
站、广场等开阔场所；吸顶扬声器则用于室内背景音乐，注重美观与均匀
覆盖。数字广播系统中，网络设备（交换机、终端解码器）需保证时钟同
步与低抖动，以避免多区声音不同步。

扬声器的匹配与功率分配直接决定系统可靠性。定压系统中每只扬声器经
线间变压器取一个功率档（如 $3/6/10/15\ \mathr
m{W}$），各档对应不同初级匝数，使高压侧等效阻抗随功率变化而
始终保持与母线匹配；多只并联后，系统等效阻抗变小、母线电流增大，
功放输出功率等于各扬声器标称功率之和。工程上应使功放额定功率略大
于扬声器总功率（预留 $20\%\sim 30\%$ 余量），并
为每区设置独立的熔断器或电子保护，防止单只短路拖累整条母线。对于
号角等高灵敏度单元，小功率即可覆盖大空间，适合广场、站台；吸顶/
壁挂单元则用于室内均匀扩声，二者常在同一系统中按场景混用。

功率放大器是有线广播的功率核心，选型除容量余量外还应关注保护功能
：过载、短路、过热与输出短路保护可避免扬声器故障扩大为功放损坏；
风机智能调速能在低负荷时降噪、高负荷时强冷。定压功放的输出变压器
兼作隔离与升压，其漏感与分布电容会影响高频响应，优质铁芯与工艺可
拓宽平坦频段。安装后应使用假负载或标准扬声器对每路输出进行满功率
老化与失真测试，确认各档抽头电压准确，并在正式接入前完成整体联调
，确保分区、强切与优先级逻辑符合设计预期。

\section{有线广播设计与安装}
\label{sec:wired_broadcast_design}

有线广播设计需在覆盖、音质与可靠性之间统筹，关键要点如下：

\begin{itemize}
    \item \textbf{功率计算}：统计各分区扬声器
数量与功率，求和后得到总负荷，功放容量取总负荷的 $1.2\si
m 1.3$ 倍。
    \item \textbf{线路设计}：依传输距离与允许
压降选择线缆规格，定压系统宜用足够线径双绞线并缩短分支长度。
    \item \textbf{分区规划}：按楼层、功能与消
防防烟分区划分广播区，确保应急时定向强切。
    \item \textbf{声学设计}：核算声场均匀度与
最大声压级，避免近处过响、远处听不清。
\end{itemize}

安装注意事项包括扬声器布局、线缆敷设、接地处理与防雷保护。扬声器
布置应保证区域内声压级差不超过数分贝；线缆宜与强电分槽敷设以减少
感应噪声；系统应一点接地并做好等电位联结；室外与高层引入端应加装
防雷与浪涌保护器。功率分配可用总功率公式校验：
\begin{equation}
P_{\mathrm{total}} = \sum_{i=1}^{N} P_i, \qquad U_0 = \sqrt{P_{\mathrm{total}} Z_{\mathrm{sys}}},
\end{equation}
其中 $Z_{\mathrm{sys}}$ 为系统折算阻抗，$U
_0$ 为所需母线电压；据此可核对功放档位与线路载流量是否匹配。

在农村应急广播体系中，通常构建"县级播控平台—乡镇分控—村终端"
三级架构：县级平台统一制作与下发预警信息，经有线（或结合无线/有
线混合）链路送达各村大喇叭与音柱，并具备断网可本地播发的兜底能力
。该类系统强调无人值守、远程巡检与多通道备份，常与气象、地震、防
汛等部门联动，是防灾减灾最后一公里的重要手段。

广播质量可用三项电声指标量化。频率响应描述系统对各频率的幅度一致
性；总谐波失真（THD）反映非线性畸变：
\begin{equation}
\mathrm{THD} = \frac{\sqrt{P_2^2+P_3^2+\cdots}}{P_1}\times 100\%,
\end{equation}
$P_1$ 为基波功率，$P_2,P_3,\dots$ 为各次谐
波功率。信噪比（SNR）则衡量有用信号相对本底噪声的大小：
\begin{equation}
\mathrm{SNR} = 10\log_{10}\frac{P_{\mathrm{signal}}}{P_{\mathrm{noise}}\ \mathrm{dB}},
\end{equation}
公共广播系统通常要求频率响应 $80\ \mathrm{Hz}\
sim 12\ \mathrm{kHz}$ 内起伏不超过数分贝、
THD 低于数个百分点、SNR 优于 $60\sim 70\ \
mathrm{dB}$，以满足清晰可懂的播报需求。

声场设计上，应利用声线法或声场模型估算各点声压级，使主要区域达到
规定的最小声压级（应急广播常要求 $\ge 65\sim 75\
 \mathrm{dB}$ 且高于背景噪声 $10\sim 15
\ \mathrm{dB}$），并控制均匀度。安装时扬声器宜均匀
布置、避免声影区；线缆与强电分槽敷设、做好屏蔽与一点接地，可有效
抑制交流哼声与感应干扰。室外与高层引入端应加装防雷与浪涌保护器，
并将金属管路、桥架与建筑接地网可靠连接。验收阶段宜实测各分区声压
级、失真与强切功能，并模拟主电中断检验备用电源（通常为蓄电池）的
续航是否满足规范要求。
